% Generado por sistema/backend ReporteTex.java -- no editar a mano.
% Datos: dataset de prueba (respuestas sinteticas)
\subsubsection*{C6. Concordancia formación–mercado $\times$ tiempo al primer empleo}

\textit{Propósito.} Determinar si percibir la formación alineada con los requerimientos del mercado se asocia con un menor tiempo para obtener el primer empleo. Respalda el criterio ARCU-SUR 3.2.3. Se aplica sobre: trabajan en una organización, con nivel de significancia $\alpha = 0.05$.

\textit{Paso 1. Contar con la muestra: tabla de contingencia} (frecuencia observada $FO_{ij}$, $N = 68$).

\begin{mitabla}{C6: frecuencia observada $FO_{ij}$ \datosprueba}
\begin{tabular}{lcc|c}
\toprule
Concordancia & Hasta 4 meses & Más de 4 meses & Total fila \\
\midrule
En desacuerdo & 5 & 13 & 18 \\
De acuerdo & 29 & 21 & 50 \\
\midrule
Total columna & 34 & 34 & $N=68$ \\
\bottomrule
\end{tabular}
\end{mitabla}

\textit{Paso 2. Calcular la frecuencia esperada} de cada celda, $FE_{ij} = (\text{total fila } i)(\text{total columna } j)/N$. Por ejemplo, para la celda ``En desacuerdo -- Hasta 4 meses'': $FE_{11} = 18 \times 34 / 68 = 9.0000$.

\begin{mitabla}{C6: frecuencia esperada $FE_{ij}$ \datosprueba}
\begin{tabular}{lcc}
\toprule
Concordancia & Hasta 4 meses & Más de 4 meses \\
\midrule
En desacuerdo & $18\cdot34/68 = 9.00$ & $18\cdot34/68 = 9.00$ \\
De acuerdo & $50\cdot34/68 = 25.00$ & $50\cdot34/68 = 25.00$ \\
\bottomrule
\end{tabular}
\end{mitabla}

\textit{Paso 3. Calcular el estadístico chi-cuadrado,} sumando $(FO_{ij}-FE_{ij})^2/FE_{ij}$ de cada celda.

\begin{mitabla}{C6: cálculo de $X_0^2$ celda por celda \datosprueba}
\begin{tabular}{L{5.4cm}cccc}
\toprule
Celda & $FO_{ij}$ & $FE_{ij}$ & $(FO_{ij}-FE_{ij})^2$ & $(FO_{ij}-FE_{ij})^2/FE_{ij}$ \\
\midrule
En desacuerdo -- Hasta 4 meses & 5 & 9.0000 & 16.0000 & 1.7778 \\
En desacuerdo -- Más de 4 meses & 13 & 9.0000 & 16.0000 & 1.7778 \\
De acuerdo -- Hasta 4 meses & 29 & 25.0000 & 16.0000 & 0.6400 \\
De acuerdo -- Más de 4 meses & 21 & 25.0000 & 16.0000 & 0.6400 \\
\midrule
\multicolumn{4}{r}{$X_0^2 = \sum$} & 4.8356 \\
\bottomrule
\end{tabular}
\end{mitabla}

\textit{Paso 4. Determinar el estadístico teórico.} Con $k = 2$ filas y $m = 2$ columnas, los grados de libertad son $(k-1)(m-1) = (2-1)(2-1) = 1$. De la tabla de chi-cuadrada: $X_{0.05,\,1}^2 = 3.841$.

\textit{Paso 5. Establecer la relación entre estadísticos:} $X_0^2 \ge X_{0.05,\,1}^2 \;\Rightarrow\; 4.8356 \ge 3.841$.

\textit{Paso 6. Conclusión.}

\begin{center}
\begin{tabular}{l}
Si ($X_0^2 < X_{0.05,\,1}^2$)\\
\quad Independencia\\
Sino\\
\quad \fbox{Relación entre las variables}\\
FinSi
\end{tabular}
\end{center}

Se rechaza la independencia: hay evidencia de que concordancia y tiempo al primer empleo están relacionadas. Proporción por fila en la primera columna (``Hasta 4 meses''): En desacuerdo $5/18 = 27.8\%$; De acuerdo $29/50 = 58.0\%$. \datosprueba

